\section{長さ二の基点による四次障害}
\label{sec:basepoint-length-two}
一点基点の一般階数問題を残したままでも、複数点と非被約基点は一回Torの部分複体から排除できる。
以下では$\dim B_1=4$, $\dim B_2=6$, $\dim B_3=4$を仮定し、
基点スキームを$Z=\operatorname{Proj}B$とする。

\begin{theorem}[長さ二の基点の排除]\label{thm:length-two-basepoint-obstruction}
$Z$が長さ二の部分スキームを含むとする。
内部次数4、位置1から位置1への一回Tor複体は
\[
 0\longrightarrow D_3\longrightarrow D_2\longrightarrow D_1
 \longrightarrow D_0\longrightarrow0,
 \qquad
 D_h=\bigoplus_{\substack{p+q=h\\0\leq p\leq2,\ 0\leq q\leq1}}
 Q_p\otimes B_h^*\otimes Q_q
\]
であり、次元は$(1,32,156,160)$である。その最上ホモロジーは42次元以上である。
従って$T_L$はKoszulではない。
さらに$w$が有限Frobenius候補なら、その元四窓について位置1に対応する始点で
$\dim C_{i,i+4}\geq42$、従って四次障害$\delta_i\geq7$となる。
特に三因子対称な標準Koszul AS候補の基点スキームは、空か、ただ一つの縮約点である。
\end{theorem}
\begin{proof}
位置と内部次数を命題\ref{prop:one-return-bar-tor}へ代入すると、上の複体を得る。
両側の最大Koszul次数は2と1であり、$M=D A_L$の成分は$B_h^*$となる。

まず一点$[v]\in Z$を取る。$F_h=k v^{\otimes h}\subset B_h^*$は縮約で保たれ、
一点評価の部分複体を与える。
その左側は$Q_2\to Q_1\to Q_0$、右側は$Q_1\to Q_0$のテンソル積である。
一次式の微分は$v$での評価、二次関係の微分は一因子の$v$縮約である。
$L$の二次式は$v$で零なので、$Q_2\to Q_1$の像は$\ker(v:V^*\to k)$に含まれる。
交換子の縮約がこの三次元空間を尽くすので、左側のホモロジーは次数2の七次元だけ、
右側のホモロジーは次数1の三次元だけである。
K\"unnethの公式により一点部分複体のホモロジーは次数3の$7\cdot3=21$次元だけとなる。

長さ二のスキームが相異なる二点なら、両点の三次cubeは独立である。
従ってその二つの21次元核は$D_3$内で交わらず、42次元の核を与える。
ここには$D_4$がないので、そのまま最上ホモロジーに入る。

長さ二のスキームが一点に支持されるなら、標数零でそれは接ベクトル$v+\epsilon u$、
$\epsilon^2=0$、$u\notin kv$で表せる。
次数$h\geq1$で二次元のジェット評価双対$F_h\subset B_h^*$を取り、$F_0=k$とする。
この部分複体は一点部分複体を含む。その商は次数0の項が零で、正次数では一次元である。
実際、$v^{\otimes h}$を除いたジェットの縮約は、商では一次式の$v$評価に等しい。
従って商複体は一点複体の次数0の項だけを除いたものであり、
次数3のホモロジーは21次元、次数2のホモロジーは零である。
部分・全体・商のホモロジー長完全列と、一点部分の次数2ホモロジーの消滅から、
ジェット部分複体の次数3ホモロジーは$21+21=42$次元となる。
これも$D_3$内の核として全複体へ単射する。

命題\ref{prop:one-return-bar-tor}により、この核は
$\operatorname{Tor}^{T_L}_{4,4}(S,S)$の指定頂点成分の直和因子である。
$T_L$の各始点のHilbert列は$(1,4,6,4,1)$なので、Koszulならその成分の次元は
$(1-t)^{-4}$の四次係数35でなければならず、矛盾する。
有限Frobeniusの場合には$E(w)=T_L$であり、二次双対の四次元数は
$\operatorname{Tor}^{E(w)}_{4,4}$の次元だから、述べた四次障害を得る。

最後に、代数閉体上の射影スキームが空でも一縮約点でもなければ、相異なる二閉点、
または一閉点の非零接方向からなる長さ二の部分スキームを含む。
従って最後の結論も従う。
\end{proof}

\begin{corollary}[残る一点枝の縮約八点表示]\label{cor:simple-basepoint-reduced-net}
三因子対称な有限Frobenius候補が全四次Koszul検査を通り、基点を持つなら、その四二次式空間$L$は
相異なる八点完全交叉を定める三二次式と、そのうち一点だけで消える第四式を基底として持つ。
従って\textup{命題\ref{prop:reduced-onepoint-commutative-betti}}の可換最小分解は、
この残余AS候補枝全体に適用できる。ただしその分解だけではAS排除を結論しない。
\end{corollary}
\begin{proof}
前定理から基点スキームは一縮約点$p$である。
$L$のイデアルの高さは3なので、一般の三式は正則列となる。
$p$では四式の一次微分が三次元余接空間を生成するため、一般の三式はそこで横断的に交わる。
$U=\mathbb P^3\setminus\{p\}$上では$L$の評価は各点で全射である。
従って三式とその共通零点のincidence
$\{(x,f_1,f_2,f_3):x\in U,\ f_i(x)=0\}\to U$は階数9のベクトル束であり、滑らかである。
パラメータ空間$L^3$の一般点で、三式は正則列かつ$p$で横断的なので、
その完全交叉の残る長さは$8-1=7$である。従ってincidenceの$L^3$への射は支配的である。
その一般ファイバーは滑らかなincidenceの局所化として正則であり、
パラメータの函数体は標数零なので、この零次元ファイバーは幾何学的に縮約である。
滑らかさの開性から一般の三式の交叉は$U$上でも縮約となる。
従って次数$2^3=8$の相異なる八点完全交叉を得る。
これは本状況でのBertiniの議論をincidenceで書いたものである。
$L$の基底をなす第四式は、他の七点で消えればそこも$L$全体の基点となるので、
$p$でのみ消える。可換最小分解の命題を適用する。
\end{proof}

\begin{remark}
\path{computations/verify_basepoint_length_two.py}は、基点自由な対照点、二点を持つ例、
接ベクトルを含む非被約例について、三次商次元4、微分の零合成と全階数を有理数で検算する。
四次の最上ホモロジーはそれぞれ35,46,84であり、普遍的下界42と整合する。
後二例が有限Frobeniusであるとは主張せず、定理の証明もこの標本の外挿ではない。
\end{remark}
